\section{Metodologia e implementazione}
\label{sec:methodology}

\subsection{Organizzazione del codice}
Il codice sta nel pacchetto \code{rl\_lab}: gli algoritmi sono funzioni pure (ambiente
in ingresso, risultati in uscita), gli script orchestrano gli esperimenti, il plotting
\`e isolato dai loop di training e gli iperparametri stanno in \code{configs.py}. Le
componenti sono \code{algorithms/tabular/} (i quattro algoritmi),
\code{algorithms/deep/} (rete Nature-CNN, replay buffer, loop DQN con i flag di
ablazione e Double DQN), \code{envs/} (discretizzatore e preprocessing Atari),
\code{experiments/runner.py} (esecuzione multi-seed: ogni esperimento passa di qui,
quindi nessun risultato pu\`o essere accidentalmente a seed singolo) e \code{utils/}
(seeding e plotting).

\begin{table}[H]
\centering
\caption{Caratteristiche degli ambienti e discretizzazione adottata.}
\label{tab:envs}
\begin{tabular}{lcccc}
\toprule
Ambiente & Stato & Azioni & Reward & Discretizzazione \\
\midrule
CartPole & 4-D cont. & 2 & $+1$/passo (denso) & 10 bin/dim $\to 10^4$ stati \\
Acrobot  & 6-D cont. & 3 & $-1$/passo (sparso) & 6 bin/dim $\to 46\,656$ stati \\
Pong     & $210{\times}160$ RGB & 6 & $\pm 1$ (punto) & --- (pixel, DQN) \\
\bottomrule
\end{tabular}
\end{table}

\subsection{Discretizzazione: le scelte che contano}
Il wrapper \code{DiscretizedEnv} clippa ogni dimensione in un intervallo configurato, la
scala e la tronca a un indice intero; il risultato \`e una tupla usata come chiave della
Q-table (state aggregation, S\&B \S9.4). Due scelte incidono sui risultati e le dichiaro.

\emph{Intervalli di clipping.} Su CartPole clippo posizione e angolo ai limiti di
\textbf{terminazione} ($\pm 2.4$ e $\pm 0.21$\,rad), non a quelli dello spazio di
osservazione ($\pm 4.8$, $\pm 0.418$): oltre le soglie di terminazione l'episodio \`e
gi\`a finito, quindi usare i limiti pi\`u larghi sprecherebbe met\`a dei bin di
posizione e angolo su stati che non possono verificarsi. La scelta raddoppia la
risoluzione utile sulle due dimensioni pi\`u importanti. Le velocit\`a non hanno un
limite ufficiale e le clippo a $\pm 3.0$, valore entro cui restano comodamente durante il
bilanciamento. Su Acrobot i limiti \emph{esistono} nel sorgente Gymnasium ($4\pi$ e $9\pi$
per le due velocit\`a angolari) e li riuso esattamente: nessun bin sprecato, nessuno
saturo.

\emph{Numero di bin.} \`E il parametro che non ho voluto assumere: la
Sez.~\ref{sec:res-acrobot-bins} lo misura direttamente, con esito opposto all'attesa.

\subsection{Quale policy valutano MC e TD(0)?}
\label{sec:met-prediction}
Addestro Q\hbox{-}Learning, ne estraggo la policy greedy $\pi(s)=\arg\max_a Q(s,a)$ e
valuto \emph{quella} con entrambi gli stimatori; come riferimento valuto anche una policy
\emph{casuale}. Il doppio caso non \`e ridondante ed \`e anzi la parte pi\`u interessante
del risultato (Sez.~\ref{sec:res-cartpole-pred}): la policy appresa genera episodi lunghi
($\approx 330$ passi), quella casuale episodi corti ($\approx 22$), e la lunghezza degli
episodi \`e esattamente ci\`o che determina quanto TD(0) resta distorto.

Per confrontare gli stimatori traccio la stima $V(s_{\text{ref}})$ di un singolo stato,
episodio per episodio. Come $s_{\text{ref}}$ uso lo \textbf{stato di reset
discretizzato}: \`e sempre raggiungibile ed \`e visitato all'inizio di ogni episodio,
quindi entrambi gli stimatori lo aggiornano di continuo. La mia prima scelta --- il
``bin centrale'' di ogni dimensione --- era sbagliata: su Acrobot d\`a uno stato
\emph{fisicamente impossibile}, perch\'e l'osservazione codifica gli angoli come
$[\cos\theta, \sin\theta]$ e i due non possono essere entrambi $\approx 0$. Quello stato
non veniva mai visitato e la stima restava piatta a zero.

\subsection{DQN, preprocessing Atari e riproducibilit\`a}
La rete \`e la Nature-CNN~\cite{mnih2015}: tre layer convoluzionali e due fully connected
($3136 \to 512 \to n_{\text{azioni}}$), che restituiscono $Q(s,\cdot)$ per tutte le
azioni in un forward pass. I due stabilizzatori sono dietro un flag ciascuno, cos\`i
l'ablazione richiesta \`e un cambio di parametro e non di codice. Il replay buffer \`e
limitato a 100k transizioni: il milione del paper Nature occuperebbe $\approx 28$\,GB,
pi\`u della RAM disponibile.

Il preprocessing Atari segue lo standard di Mnih~et~al.: no-op iniziali casuali,
frame-skip 4 con massimo pixel-wise sugli ultimi due frame (Atari fa lampeggiare gli
sprite e la pallina pu\`o sparire a frame alterni), grayscale, resize $84{\times}84$ e
stack di 4 frame --- necessario perch\'e da un solo fotogramma la direzione della pallina
\`e ambigua. I frame restano \code{uint8} nel buffer e sono normalizzati solo al momento
del training.

Ogni run fissa il seed con \code{set\_seed} (numpy + torch) e
\code{env.reset(seed=\ldots)}; tutti i risultati sono media su 3 seed $\{42, 123, 456\}$.
Gli array grezzi sono salvati in \code{results/} come \code{.npz} e \textbf{tutte le
tabelle di questo report sono generate da quei file}
(\code{report/make\_report\_tables.py}): nessun numero \`e trascritto a mano.

\subsection{Origine del codice}
\label{sec:code-origin}
I \textbf{quattro algoritmi tabellari} sono adattati dalle esercitazioni del corso
(\code{control.py}, \code{prediction.py}): stessa regola di aggiornamento, riadattata
all'API Gymnasium moderna e agli stati discretizzati. Il \textbf{DQN} (replay buffer,
loop di training, target network, Double DQN), il \textbf{discretizzatore},
l'infrastruttura multi-seed e tutti gli script sono \textbf{scritti da zero}: il corso
forniva solo un DQN basato su Stable-Baselines3, che la consegna vieta per il core. La
Nature-CNN replica l'architettura di~\cite{mnih2015} con implementazione PyTorch
originale. Nessuna riga di Stable-Baselines3 \`e usata nel progetto. Il dettaglio
completo \`e in \code{docs/origine\_codice.md}.
